\pagenumbering{gobble}

\selectlanguage{magyar}
\hungarianParagraph

%----------------------------------------------------------------------------
% Abstract in Hungarian
%----------------------------------------------------------------------------

\chapter*{Kivonat}

A dolgozatom egy mesterséges intelligenciával támogatott beszédfejlesztő alkalmazás tervezését és megvalósítását mutatja be. A hagyományos nyelvtanulási módszerek gyakran a nyelvtani szabályok és az elszigetelt szókincs elsajátítására helyezik a hangsúlyt, miközben korlátozott lehetőséget biztosítanak a valós kommunikáció gyakorlására. Ennek következtében a tanulók nehezebben alkalmazzák tudásukat valódi beszédhelyzetben.

A fejlesztett alkalmazás célja, hogy beszélgetésalapú tanulási környezetet biztosítson, ahol a felhasználók mesterséges intelligenciával folytatott párbeszédek során fejleszthetik nyelvi készségeiket. Az AI automatikusan felismeri, kijavítja a hibáikat, majd segít nekik a fejlődésben azzal hogy ismételten előkerülnek ezek a beszélgetések során.

A dolgozat bemutatja az alkalmazás felépítését, a megvalósítás során alkalmazott módszereket és technológiai megoldásokat, valamint azt, hogy a beszélgetésközpontú, hibaalapú tanulási megközelítés miként járult hozzá a beszédkészség hatékony fejlesztéséhez.

\vfill
\selectlanguage{romanian}

%----------------------------------------------------------------------------
% Abstract in Romanian
%----------------------------------------------------------------------------
\chapter*{Rezumat}

Lucrarea mea prezintă proiectarea și implementarea unei aplicații de dezvoltare a abilităților de vorbire, susținută de inteligență artificială. Metodele tradiționale de învățare a limbilor străine pun adesea accentul pe însușirea regulilor gramaticale și a vocabularului izolat, oferind în același timp oportunități limitate pentru exersarea comunicării reale. Ca urmare, cursanții întâmpină dificultăți în aplicarea cunoștințelor dobândite în situații reale de vorbire.

Scopul aplicației dezvoltate este de a oferi un mediu de învățare bazat pe conversație, în care utilizatorii își pot dezvolta competențele lingvistice prin dialoguri purtate cu o inteligență artificială. AI-ul identifică și corectează automat greșelile utilizatorilor și sprijină procesul de învățare prin reluarea acestora în cadrul conversațiilor ulterioare.

Lucrarea prezintă structura aplicației, metodele și soluțiile tehnologice utilizate în procesul de implementare, precum și modul în care abordarea de învățare centrată pe conversație și pe greșeli contribuie la dezvoltarea eficientă a abilităților de comunicare orală.

\vfill
\selectlanguage{english}
%\englishParagraph

%----------------------------------------------------------------------------
% Abstract in English
%----------------------------------------------------------------------------
\chapter*{Abstract}

This thesis presents the design and implementation of an artificial intelligence–supported application for the development of speaking skills. Traditional language learning methods often emphasize the acquisition of grammatical rules and isolated vocabulary, while providing limited opportunities for practicing real-life communication. As a result, learners face difficulties in applying their knowledge in authentic speaking situations.

The aim of the developed application is to provide a conversation-based learning environment in which users can improve their language skills through dialogues with an artificial intelligence. The AI automatically identifies and corrects users’ errors and supports their progress by revisiting these mistakes during subsequent conversations.

The thesis describes the structure of the application, the methods and technological solutions applied during its implementation, and examines how a conversation-centered, error-based learning approach contributes to the effective development of speaking skills.


\vfill
\dolgozatnyelve
\defaultParagraph